% Options for packages loaded elsewhere
\PassOptionsToPackage{unicode}{hyperref}
\PassOptionsToPackage{hyphens}{url}
\PassOptionsToPackage{dvipsnames,svgnames,x11names}{xcolor}
\documentclass[
  ngerman,
  10pt,
  a4paper,
]{article}
\usepackage{xcolor}
\usepackage[margin=22mm]{geometry}
\usepackage{amsmath,amssymb}
\setcounter{secnumdepth}{-\maxdimen} % remove section numbering
\usepackage{iftex}
\ifPDFTeX
  \usepackage[T1]{fontenc}
  \usepackage[utf8]{inputenc}
  \usepackage{textcomp} % provide euro and other symbols
\else % if luatex or xetex
  \usepackage{unicode-math} % this also loads fontspec
  \defaultfontfeatures{Scale=MatchLowercase}
  \defaultfontfeatures[\rmfamily]{Ligatures=TeX,Scale=1}
\fi
\usepackage{lmodern}
\ifPDFTeX\else
  % xetex/luatex font selection
  \setmainfont[Path=/Users/stefanhamann/.cache/codex-runtimes/codex-primary-runtime/dependencies/native/libreoffice-headless/libreoffice/LibreOfficeDev.app/Contents/Resources/fonts/truetype/,Extension=.ttf,BoldFont=DejaVuSans-Bold,ItalicFont=DejaVuSans-Oblique,BoldItalicFont=DejaVuSans-BoldOblique]{DejaVuSans}
  \setmonofont[Path=/Users/stefanhamann/.cache/codex-runtimes/codex-primary-runtime/dependencies/native/libreoffice-headless/libreoffice/LibreOfficeDev.app/Contents/Resources/fonts/truetype/,Extension=.ttf,BoldFont=DejaVuSansMono-Bold,ItalicFont=DejaVuSansMono-Oblique,BoldItalicFont=DejaVuSansMono-BoldOblique]{DejaVuSansMono}
\fi
% Use upquote if available, for straight quotes in verbatim environments
\IfFileExists{upquote.sty}{\usepackage{upquote}}{}
\IfFileExists{microtype.sty}{% use microtype if available
  \usepackage[]{microtype}
  \UseMicrotypeSet[protrusion]{basicmath} % disable protrusion for tt fonts
}{}
\makeatletter
\@ifundefined{KOMAClassName}{% if non-KOMA class
  \IfFileExists{parskip.sty}{%
    \usepackage{parskip}
  }{% else
    \setlength{\parindent}{0pt}
    \setlength{\parskip}{6pt plus 2pt minus 1pt}}
}{% if KOMA class
  \KOMAoptions{parskip=half}}
\makeatother
\ifLuaTeX
\usepackage[bidi=basic,shorthands=off]{babel}
\else
\usepackage[bidi=default,shorthands=off]{babel}
\fi
\ifPDFTeX
\else
\babelfont{rm}[Path=/Users/stefanhamann/.cache/codex-runtimes/codex-primary-runtime/dependencies/native/libreoffice-headless/libreoffice/LibreOfficeDev.app/Contents/Resources/fonts/truetype/,Extension=.ttf,BoldFont=DejaVuSans-Bold,ItalicFont=DejaVuSans-Oblique,BoldItalicFont=DejaVuSans-BoldOblique]{DejaVuSans}
\fi
\ifLuaTeX
  \usepackage{selnolig} % disable illegal ligatures
\fi
\setlength{\emergencystretch}{3em} % prevent overfull lines
\providecommand{\tightlist}{%
  \setlength{\itemsep}{0pt}\setlength{\parskip}{0pt}}
\usepackage{fvextra}
\RecustomVerbatimEnvironment{verbatim}{Verbatim}{breaklines=true,fontsize=\footnotesize}
\usepackage{xurl}
\usepackage{seqsplit}
\usepackage{adjustbox}
\usepackage{ragged2e}
\usepackage{titlesec}
\titleformat{\section}{\large\bfseries\raggedright}{}{0em}{}
\titleformat{\subsection}{\normalsize\bfseries\raggedright}{}{0em}{}
\DeclareRobustCommand{\codeword}[1]{\texttt{\seqsplit{#1}}}
\AtBeginEnvironment{longtable}{\small\RaggedRight}
\usepackage{newunicodechar}
\newunicodechar{𝒞}{\ensuremath{\mathcal C}}
\newunicodechar{𝔊}{\ensuremath{\mathfrak G}}
\newunicodechar{𝔤}{\ensuremath{\mathfrak g}}
\newunicodechar{𝓗}{\ensuremath{\mathcal H}}
\newunicodechar{ℋ}{\ensuremath{\mathcal H}}
\newunicodechar{𝔄}{\ensuremath{\mathfrak A}}
\newunicodechar{𝟙}{\ensuremath{\mathbf 1}}
\newunicodechar{𝕀}{\ensuremath{I}}
\newunicodechar{𝕋}{\ensuremath{\mathbb T}}
\newunicodechar{𝟘}{\ensuremath{\mathbf 0}}
\newunicodechar{⊞}{\ensuremath{\boxplus}}
\newunicodechar{∘}{\ensuremath{\circ}}
\newunicodechar{⟦}{\ensuremath{[\![}}
\newunicodechar{⟧}{\ensuremath{]\!]}}
\setlength{\emergencystretch}{4em}
\setlength{\parindent}{0pt}
\setlength{\parskip}{0.45em}
\AtBeginDocument{\hypersetup{colorlinks=true,linkcolor=black,urlcolor=blue}\pdfstringdefDisableCommands{\def\codeword#1{#1}}\RaggedRight}
\usepackage{fancyhdr}
\pagestyle{fancy}
\fancyhf{}
\fancyhead[L]{\small TFPT / Universalraum - v1.6.10}
\fancyfoot[R]{\thepage}
\setlength{\headheight}{15pt}
% The preserved full history now has three-digit page numbers.
\makeatletter
\renewcommand{\@pnumwidth}{2.6em}
\renewcommand{\@tocrmarg}{3.6em}
\makeatother

\usepackage{bookmark}
\IfFileExists{xurl.sty}{\usepackage{xurl}}{} % add URL line breaks if available
\urlstyle{same}
\hypersetup{
  pdflang={de-DE},
  colorlinks=true,
  linkcolor={Maroon},
  filecolor={Maroon},
  citecolor={Blue},
  urlcolor={Blue},
  pdfcreator={LaTeX via pandoc}}

\author{}
\date{}

\begin{document}

\section{TFPT / Universalraum - Update
v1.6.10}\label{tfpt-universalraum---update-v1.6.10}

\begin{enumerate}
\def\labelenumi{\arabic{enumi}.}
\setcounter{enumi}{14}
\tightlist
\item
  September 2026. Bezug: Hauptfassung v1.6.9 und die neu geprüfte Arbeit
  „Gemeinsame Quelle``, Fortsetzung 1.2.
\end{enumerate}

\subsection{Die neue Verbindung}\label{die-neue-verbindung}

TFPT enthält bereits einen fermionischen Vierteldrehungsclock mit
\(\rho^4=-I\). Das neue Quellpaper zeigt gleichzeitig, dass
Paarantworten ein solches Vorzeichen verlieren können. Diese beiden
Befunde wurden nun an derselben festgelegten Viererquelle mit dem
unveränderten inneren Tensor W zusammengeführt.

Die Quelle verwendet \(q=(I+R_\eta)f/\sqrt2\), \(R_\eta^4=\eta I\), vier
Fermion- und vier Bosonbereiche. Das ist eine ausdrückliche
Kompositionsannahme, noch keine aus TFPT ausgewählte Raumgeometrie. Die
zusätzliche Kopienstruktur ist nicht die innere SU(4)-Familienstruktur.

\subsection{Was neu bewiesen wurde}\label{was-neu-bewiesen-wurde}

\begin{itemize}
\tightlist
\item
  Die periodische Variante besitzt 64 völlig ungekoppelte Fermionmoden;
  die antiperiodische besitzt keine. Ihre gesamten eingeschränkten
  bosonischen N=2-Antworten sind dennoch gleich.
\item
  Ein zusätzlicher Fermion-Eingang macht den Unterschied ladungsneutral
  sichtbar. Für denselben gemischten Einboson-Einfermion-Eingang
  unterscheiden sich die Bosonrückkehrwahrscheinlichkeiten erstmals
  durch \(-g^8t^8/168+O(t^{10})\).
\item
  Für \(0<|g|/\Delta\le1/2000\) hat die antiperiodische Quelle einen
  \textbf{eindeutigen Grundzustand des vollständigen Hamiltonoperators}
  bei Gesamtladung 256, ohne \(\mu N\)-Zusatz und ohne Bosonabschneiden.
  Die Lücke ist größer als \(g^2/\Delta\). Der Beweis benötigt nur die
  erneut verifizierte Casimir-Identität, Ladungserhaltung und eine
  Zweizustands-Ritzschranke.
\item
  Für den engeren Bereich \(0<|g|/\Delta\le1/10000\) ist die Antwort
  ursprünglicher clockaufgelöster Fermionen auf genau diesem
  Grundzustand kontrolliert: zwei getrennte niedrige Entnahmelinien,
  jeweils Vielfachheit 128. Die zu einem Clock-Eigenmodus gehörende
  Linie trägt über 99.992 Prozent des normierten Entnahmespektrums.
\end{itemize}

Die Energien liegen in Einheiten \(g^2/\Delta\) in den getrennten
Intervallen

\[\adjustbox{max width=\linewidth}{$\displaystyle \epsilon_-:\ (3.8933,4.6435),\qquad \epsilon_+:\ (25.1066,25.8567).$}\]

Das sind nach außen gerundete strenge Schranken aus einem summierten
Operatorrest, keine berechneten Zentralwerte. Ein einzelner
ursprünglicher Kopienmodus enthält beide Clock-Paare und sieht im
Allgemeinen beide Linien.

\subsection{Was nicht verwechselt werden
darf}\label{was-nicht-verwechselt-werden-darf}

Der bisherige Prüfpunkt \(g/\Delta=1/20\) liegt außerhalb beider neuen
Bereiche. Der neue N=256-Grundzustand ist nicht der alte
N=64-Grundzustand einer einzelnen Bank. Der N=3-Vergleich ist ein
Diagnoseexperiment und nicht die Grundzustandsantwort. Die C8-Hebung
wird nicht mit dem inneren Perioden-sechs-Clock gleichgesetzt.

Die ursprüngliche Quellen-/Kopienregel, die physische Parameterauswahl,
ein nativer Aktuator und ein beschreibbares Record bleiben offen. Ebenso
bleiben sämtliche vollständigen T1-T8-Gates offen, insbesondere
3+1D-Raumzeit, chirales Maß und dynamischer Spin 2.

\subsection{Zwei methodische
Korrekturen}\label{zwei-methodische-korrekturen}

Die volle assoziative Erzeugung durch quadratische Cliffordwörter ist
nicht volle dynamische Steuerbarkeit. Ein einzelner zusätzlicher
quartischer Generator würde unter gewährten quadratischen Kontrollen
genügen; seine native Realisierung ist nicht bewiesen.

Endlich verschiedene Mikrographen müssen nicht verschiedene großskalige
Universalitätsklassen ergeben. Deshalb sollte die Suche nicht an
vollständiger mikroskopischer Eindeutigkeit hängen bleiben. Die alte
logarithmische E8-Kaskade ordnet Skalen; sie ist nicht schon eine
physische Zeitentwicklung. Ihr archiviertes \(\gamma_0\) bleibt auch in
Ratios über den Exponenten wirksam.

\subsection{Die beiden zuletzt eingegangenen
Perspektiven}\label{die-beiden-zuletzt-eingegangenen-perspektiven}

Der vollständige positive Prozesskern ist ein tragfähiger
Rekonstruktionsrahmen, bestimmt aber seine eigenen Daten und Instrumente
nicht. Positivität allein genügt ohne passende Multiplikationsrelationen
nicht. Gleicher Grundzustand und gleiche statische Beobachtungen
bestimmen zudem keine eindeutigen Anregungsenergien.

Die Idee einer aus dem Zustand gewonnenen modularen Zeit hat einen
präzisen Test: Auf einer vollen endlichen Matrixalgebra stimmt sie genau
dann mit der vorhandenen Hamiltonentwicklung bis auf Zeitskalierung
überein, wenn der unabhängig gewonnene Zustand Gibbsform für diesen H
hat. Aus H erst den Gibbszustand zu bilden wäre keine neue Herleitung.
Der reine globale Grundzustand ist auf der vollen Algebra nicht treu;
erst eine begründete Teilalgebra kann eine geeignete modulare
Beschreibung liefern.

Neu abgegrenzt wurde außerdem die Kausalitätsspur: Eine Automorphie
erhält die Dimension einer endlichen Algebra. Eine halbseitige Inklusion
ihres Bildes in sich selbst muss dort Gleichheit sein. Nichttriviale
modulare Raumzeitrekonstruktion benötigt somit mehr als einen endlichen
Clockblock, etwa einen kontrollierten unendlichen Grenzübergang.

Die Hinweise auf die Gesamtsymmetrie und verborgene Schleifenphasen
tragen. Unser Kopienclock bricht allerdings die innere Symmetrie nicht.
Das allgemeine Gegenargument mit höheren Bosonzahltermen trägt
ebenfalls: Eine endliche Zahl erhaltener Testsektoren wählt keinen
unbegrenzten Hamiltonoperator aus. Die 61 zusätzlichen exakten
Prüfbedingungen liefen normal und optimiert mit identischen Resultaten;
die allgemeinen Argumente stehen in Abschnitt 14 des Hauptdokuments.

\subsection{Weitergerechnet: Auswahl durch gemeinsame
Spiegelung?}\label{weitergerechnet-auswahl-durch-gemeinsame-spiegelung}

Clock und Grundzustand wählen die Verdrehung allein \textbf{nicht}: Die
reelle Quelle \(q=(3I+4R_+)f/5\) hat ebenfalls keinen freien
Fermionbereich und bei genügend kleiner Kopplung einen eindeutigen
globalen Grundzustand. Sie besitzt allerdings eine andere Paarmatrix als
die fixierte gleichstarke Mischung.

Eine zusätzliche, genau erklärte gemeinsame Links-rechts-Spiegelung von
Fermionorten und Paarbänken erzwingt dagegen gleiche Gewichtsbeträge.
Dann lässt der Ausschluss exakt freier Fermionrichtungen innerhalb
dieser Zweinachbar-Klasse nur den negativen Lift zu. Die gesamte Auswahl
hängt so an einer kleinen gemeinsamen Operatorrelation, nicht an einer
zusätzlichen großen Matrix.

TFPT enthält die verwandte geometrische Relation
\(\sigma\rho\sigma=\rho^{-1}\) bereits. Noch fehlt der Nachweis ihrer
passenden gemeinsamen Wirkung auf die ursprünglichen Fermion- und
W-Paarressourcen. Abstrakte dihedrale Symmetrie oder
Reflexionspositivität allein liefern diesen Operatoradapter nicht. Der
neue bedingte Auswahlsatz samt Gegenmodell steht in Abschnitt 15.

\subsection{Priorität}\label{priorituxe4t}

Zuerst die gemeinsame Spiegelungswirkung der TFPT-Seam auf Fermionorte
und Paarübergänge prüfen; daran den zustands- und operatorverträglichen
Quellenanschluss entscheiden. Danach den Kopplungsbereich erweitern und
dieselbe geladene Antwort in einer Größenfolge untersuchen. Das ist eine
kleinere, konkrete Forschungsaufgabe als die Suche nach immer weiteren
passenden Zahlen.

Verifikation: Das externe Paket bestand erneut 614 Bedingungen normal
und optimiert. Die eigene Rechnung ist separat mit exakten CAR-,
Casimir-, Momenten- und Restschranken kontrolliert; die allgemeinen
Beweise stehen vollständig im Haupttext. Die ältere Gesamtsuite wurde
nicht noch einmal vollständig ausgeführt.

\end{document}
